% !TeX root = ../../Book3.tex
% 纲要标题：### 8.2 Buchberger 理论
% 纲要位置：工作记录/计划纲要.md，第 2148 行。
% 纲要细目（写作范围）：
% * 从首项对齐构造 S 多项式，以最大首单项式最小的表示和有界表示完整证明 Buchberger 判据
% * 给出 Buchberger 算法及完整终止性、正确性证明，以 Dickson 有限性控制首项理想升链
% * 构造约化 Gröbner 基，并完整证明固定基域和单项式序下的唯一性
% * 用手算例题与零余式表示证书解释基础算法，完整书写当前应用的推导

\section{Buchberger 理论}
\label{b3:sec:0802}

Gröbner 基控制理想中全部多项式的首项，定义本身包含无限多个检验。有限算法的关键在于理解最高项怎样抵消：两个生成元的首项抵消由一个 S 多项式记录，而多个最高项同时抵消可以分解为这样的两两抵消。本节将这项观察写成首项下降的严格论证，再据此建立算法。仍固定 \(R=k[x_1,\ldots,x_n]\) 及全局单项式序。
\subsection{S 多项式、Buchberger 判据与算法正确性}
\label{b3:subsec:0802:01}

在 \(k[x,y]\) 的字典序 \(x\succ y\) 下，取 \(f=x^2-y\)、\(g=xy-1\)。首单项式 \(x^2\) 与 \(xy\) 的最小公倍式为 \(x^2y\)，所以分别乘以 \(y\) 与 \(x\)，就能使两个首项相同，再相减得到
\[
yf-xg=y(x^2-y)-x(xy-1)=x-y^2.
\]
原来的最高单项式 \(x^2y\) 已经消失。对于非首一多项式，还需除以各自的首系数，使这两个最高项的系数也相同。

\begin{definition}\label{b3:def:s-polynomial}
设 \(k\) 为域，\(R=k[x_1,\ldots,x_n]\)，并固定全局单项式序。对非零 \(f,g\in R\)，置
\[
L=\operatorname{lcm}\bigl(\operatorname{LM}(f),\operatorname{LM}(g)\bigr).
\]
称
\begin{equation}\label{b3:eq:s-polynomial}
S(f,g)=
\frac{L}{\operatorname{LM}(f)}\frac{f}{\operatorname{LC}(f)}
-\frac{L}{\operatorname{LM}(g)}\frac{g}{\operatorname{LC}(g)}
\end{equation}
为 \(f,g\) 的 \term[S-duoxiangshi]{\(S\) 多项式}[S-polynomial]。单项式最小公倍式的各个指数取两指数的较大者。
\end{definition}
\symidx{\(S(f,g)\)：多项式 \(f,g\) 的 S-多项式}

两项的首项都归一化成 \(L\)，故它们相减后或者为零，或者首单项式严格小于 \(L\)。仅有 \(S(f,g)\in(f,g)\) 没有检验意义；重要的是能否用原生成元把它表示出来，同时让表示中的每个乘积也严格小于 \(L\)。

\begin{theorem}[Buchberger 判据]\label{b3:thm:buchberger-criterion}
设 \(k\) 为域，\(R=k[x_1,\ldots,x_n]\)，并固定全局单项式序。设 \(G=(g_1,\ldots,g_s)\) 由 \(R\) 中的非零多项式组成，\(I=(G)\)。以下条件等价：
\begin{equivalences}
\item\label{b3:item:buchberger-basis} \(G\) 是 \(I\) 的 Gröbner 基；
\item\label{b3:item:buchberger-zero} 每个 \(S(g_i,g_j)\)、\(i<j\)，按 \(G\) 除法的余式为零；
\item\label{b3:item:buchberger-bound} 对每个 \(i<j\)，存在表示
\begin{equation}\label{b3:eq:s-bounded-representation}
S(g_i,g_j)=\sum_{\ell=1}^s h_{ij\ell}g_\ell,
\qquad
\operatorname{LM}(h_{ij\ell}g_\ell)\prec L_{ij}
\quad(h_{ij\ell}\ne0),
\end{equation}
其中 \(L_{ij}=\operatorname{lcm}\bigl(\operatorname{LM}(g_i),\operatorname{LM}(g_j)\bigr)\)。
\end{equivalences}
空组对应零理想，后两个条件此时为空条件。
\end{theorem}
\begin{proof}
若\itemref{b3:item:buchberger-basis}成立，S 多项式属于 \(I\)，由\cref{b3:thm:groebner-existence-normal-form}余式为零，得到\itemref{b3:item:buchberger-zero}。再用多元除法的首项界，每个非零乘积的首单项式不超过 S 多项式的首单项式，而后者严格小于 \(L_{ij}\)，便得\itemref{b3:item:buchberger-bound}。S 多项式为零时取所有系数为零。

反过来，假定有这些有界表示。给定 \(0\ne f\in I\)，要证明某个 \(\operatorname{LM}(g_i)\) 整除 \(\operatorname{LM}(f)\)，我们希望找到一个生成元表示，使最大的乘积首单项式恰为 \(\operatorname{LM}(f)\)。若这个最大首单项式比 \(\operatorname{LM}(f)\) 更大，它的系数必在求和时抵消；S 多项式的有界表示正用来重新组织这些抵消。记 \(\widetilde g_i=g_i/\operatorname{LC}(g_i)\)，将生成元表示的系数展开为项，可写成
\[
f=\sum_{\nu=1}^t a_\nu m_\nu\widetilde g_{i_\nu},
\qquad a_\nu\in k^\times,\quad m_\nu\;\text{为单项式}.
\]
在全部这样的表示中，选择
\[
L_{\max}=\max_\nu\bigl\{m_\nu\operatorname{LM}(g_{i_\nu})\bigr\}
\]
最小的一种，良序性保证它存在。相加不会产生更大的项，故 \(\operatorname{LM}(f)\preceq L_{\max}\)。如果相等，某个 \(\operatorname{LM}(g_i)\) 整除 \(\operatorname{LM}(f)\)，已经得到所需结论。

否则令 \(J\) 为达到 \(L_{\max}\) 的那些指标。最高项系数抵消给出 \(\sum_{\nu\in J}a_\nu=0\)。固定 \(\nu_0\in J\)，于是
\[
\sum_{\nu\in J}a_\nu m_\nu\widetilde g_{i_\nu}
=\sum_{\nu\in J\setminus\{\nu_0\}}
a_\nu\bigl(m_\nu\widetilde g_{i_\nu}
-m_{\nu_0}\widetilde g_{i_{\nu_0}}\bigr).
\]
括号中若两个生成元下标相同，则乘子也相同，差为零；否则，令它们首单项式的最小公倍式为 \(L\)。二者都整除 \(L_{\max}\)，故 \(L\mid L_{\max}\)，括号中的差恰为 \((L_{\max}/L)S(g_{i_\nu},g_{i_{\nu_0}})\)。下标反序只改变符号。按假设，这可以用各 \(g_\ell\) 表示，且每个非零乘积的首单项式严格小于 \(L_{\max}\)。系数展开成项后这个界仍保持。

替换全部达到 \(L_{\max}\) 的项，其他原项本来就小于 \(L_{\max}\)，便得到最大首单项式更小的表示，与所选 \(L_{\max}\) 的最小性矛盾。因此必有 \(\operatorname{LM}(f)=L_{\max}\)。每个非零理想成员的首单项式都被某个 \(\operatorname{LM}(g_i)\) 整除，结合 \(G\subset I\) 即得首项理想相等，证明\itemref{b3:item:buchberger-basis}。
\end{proof}

这称为\term[Buchberger-panju]{Buchberger 判据}[Buchberger criterion]。证明中的第三种表述不仅解释正确性，还使我们能够保存每次计算的具体依据：以后即使加入新生成元，旧的有界表示仍然成立，不必假定不同除法顺序给出相同的中间余式。

\subsection{首项理想升链与算法终止性}
\label{b3:subsec:0802:02}

如果一个 S 多项式留下非零余式，余式仍在原理想中，但它的首单项式未被当前生成元控制。把余式加入即可改进生成组；同时必须登记它与原有全部生成元组成的新多项式对。

\begin{algorithm}[htbp]
\caption{Buchberger}
\label{b3:alg:buchberger}
\begin{algorithmsteps}{有限组 \(F=(f_1,\ldots,f_m)\subseteq R\)，可有效比较的全局单项式序，以及支持精确四则运算和零判定的系数域 \(k\)。}{\(I=(F)\) 的一个有限 Gröbner 基 \(G\)。}
\State 删除零多项式，将其余元素首一化，记为 \(G=(g_1,\ldots,g_s)\)。
\State 建立待处理集合 \(P=\bigl\{(i,j)\mid1\le i<j\le s\bigr\}\)。
\While{\(P\ne\varnothing\)}
\State 取出并删除一对 \((i,j)\)，按当前 \(G\) 对 \(S(g_i,g_j)\) 除法，得到余式 \(r\)。
\If{\(r\ne0\)}
\State 把 \(r/\operatorname{LC}(r)\) 加到 \(G\) 末尾，记新下标为 \(s+1\)。
\State 将 \((1,s+1),\ldots,(s,s+1)\) 加入 \(P\)，再令 \(s\) 增加一。
\EndIf
\EndWhile
\State \Return \(G\)。
\end{algorithmsteps}
\end{algorithm}

\begin{theorem}\label{b3:thm:buchberger-algorithm}
设域 \(k\) 支持精确四则运算与零判定，输入为 \(k[x_1,\ldots,x_n]\) 中的有限生成组及一个可有效比较的全局单项式序。\Cref{b3:alg:buchberger}在有限步后终止，输出确为原理想的 Gröbner 基。算法还可同时记录每个新生成元关于原输入的多项式线性组合。
\end{theorem}
\begin{proof}
每轮中的多元除法由\cref{b3:thm:multivariate-division}终止，依据是待处理首单项式在全局良序中严格下降。该轮有等式 \(S(g_i,g_j)=\sum_\ell q_\ell g_\ell+r\)，所以余式仍在当前理想中，添加它不改变理想。非零余式的首单项式不被已有首单项式整除，故每次添加使它们生成的单项式理想严格增大；\cref{b3:lem:dickson}的升链稳定性排除无限次添加。这控制了生成组的增长。每次添加只引入有限个新对，每轮删除一个待处理对；不再添加之后，待处理集合已经有限，余下的对也会在有限轮内处理完，因此整个算法终止。

处理某对时，若余式为零，多元除法给出\eqref{b3:eq:s-bounded-representation}。若余式非零，添加首一 \(g_{\rm new}=r/\operatorname{LC}(r)\) 后，有
\[
S(g_i,g_j)=\sum_\ell q_\ell g_\ell+
\operatorname{LC}(r)g_{\rm new}.
\]
所有非零乘积的首单项式都不超过原 S 多项式的首单项式，因而严格小于 \(L_{ij}\)。算法保留已有生成元，这些有界表示在后续步骤中仍成立。终止时最终组的每一对都已处理，由\cref{b3:thm:buchberger-criterion}即得 Gröbner 性。

最初首一化仅除以非零常数。已有生成元关于原输入的表示时，将它们代入 S 多项式和除法等式，就得到余式的表示，再除以首系数得到新生成元的表示。因此可以与算法同步，逐步记录全部关系。
\end{proof}

还可以省去部分 S 多项式的实际除法。例如首一 \(f=a+u,g=b+v\)，其中 \(a,b\) 为首单项式；若 \(a,b\) 互素，则
\[
S(f,g)=bf-ag=ug-vf.
\]
这里两个非零乘积的首单项式都严格小于 \(ab\)。因此这对已经具有判据所需的有界表示。这个\term[Buchberger-chengji-panju]{Buchberger 乘积判据}[Buchberger product criterion]只略去有数学依据的对；单凭两个多项式看起来简单，不能跳过检验。

\subsection{约化 Gröbner 基的唯一性}
\label{b3:subsec:0802:03}

\begin{definition}\label{b3:def:reduced-groebner}
设 \(k\) 为域，\(R=k[x_1,\ldots,x_n]\)，并固定全局单项式序。给定理想 \(I\subseteq R\) 的 Gröbner 基 \(G\)。若 \(G\) 的各元素首一，任意两个不同元素的首单项式互不整除，而且每个元素除首单项式以外的全部单项式都是 \(I\) 的标准单项式，则称 \(G\) 为\term[yuehua-Groebner-ji]{约化 Gröbner 基}[reduced Gröbner basis]。零理想的约化基为空，单位理想的约化基为 \(\{1\}\)。
\end{definition}

这两项限制分别去除首项生成组中的整除冗余，并把每个基元素的尾部完全化为标准单项式。

\begin{theorem}\label{b3:thm:reduced-groebner-unique}
设 \(k\) 为域，并在 \(R=k[x_1,\ldots,x_n]\) 上固定全局单项式序。每个理想 \(I\subseteq R\) 都有唯一的约化 Gröbner 基。由任意 Gröbner 基可以在有限步内得到它。
\end{theorem}
\begin{proof}
先在一组 Gröbner 基的首单项式中只保留整除关系下的极小者，相同单项式只保留一个对应多项式。保留者仍属于原理想，且首单项式生成原首项理想，因而仍为 Gröbner 基。将它们首一化，记为 \(G=\{g_1,\ldots,g_s\}\)，并置 \(m_i=\operatorname{LM}(g_i)=\operatorname{LT}(g_i)\)。对每个 \(i\)，用 \(G\setminus\{g_i\}\) 对整个尾部 \(g_i-m_i\) 作完整的多元除法，得到
\[
g_i-m_i=\sum_{j\ne i}q_{ij}g_j+r_i.
\]
以 \(g_i'=m_i+r_i=g_i-\sum_{j\ne i}q_{ij}g_j\) 替换 \(g_i\)。除法的首项界保证 \(r_i\) 的每个单项式都小于 \(m_i\)，所以首项不变；余式条件保证这些单项式不被任何 \(m_j\)（\(j\ne i\)）整除。它们也不能被 \(m_i\) 整除，因为 \(m_i\) 的真倍数严格大于 \(m_i\)。全部 \(g_i'\) 仍在 \(I\) 中，且其首单项式仍生成原首项理想，故它们构成约化 Gröbner 基。空基的情形无需除法。

两个约化基的首单项式都恰为首项理想中整除关系的极小单项式。这组单项式只由理想决定：任一理想单项式被其中一个整除，而极小者不能被另一个严格更小的理想单项式整除。按相同首单项式配对两个基元素 \(g,h\)。二者首一，故首项相消；所有余下项均为标准单项式。又 \(g-h\in I\)，由\cref{b3:thm:groebner-existence-normal-form}只能有 \(g-h=0\)。逐对相等，证明唯一性。
\end{proof}

唯一的是固定次序下的约化基，而不是任意 Gröbner 基。增加一个已在理想中的冗余多项式通常不会破坏 Gröbner 基性质；改变次序却可能同时改变首项理想、标准单项式和约化基。两种变化应分别考察。

\subsection{手算例题与零余式的线性组合表示}
\label{b3:subsec:0802:04}

在 \(\mathbb Q[x,y]\) 中取字典序 \(x\succ y\)，以
\[
f_1=x^2-y,\qquad f_2=xy-1
\]
为输入。前面的首项对齐已经给出 \(h=yf_1-xf_2=x-y^2\)。\(h\) 的首单项式 \(x\) 不被 \(x^2,xy\) 整除，所以要加入。随后 \(f_2\) 与 \(h\) 的 S 多项式为
\[
p=f_2-yh=y^3-1.
\]
这又提供未被控制的纯 \(y\) 首单项式。我们可以不再保留冗余的原生成元，而用以下双向等式核验最终理想：
\begin{equation}\label{b3:eq:groebner-cubic-certificates}
\begin{aligned}
h&=yf_1-xf_2,&
p&=-y^2f_1+(1+xy)f_2,\\
f_1&=(x+y^2)h+yp,&
f_2&=yh+p.
\end{aligned}
\end{equation}
因此 \(I=(f_1,f_2)=(h,p)\)。对最后一组，只需检验一对：
\[
S(h,p)=y^3h-xp=x-y^5=h-y^2p.
\]
先减去 \(h\)，再消去 \(-y^2p\)，余式为零。判据表明 \(\{x-y^2,y^3-1\}\) 是 Gröbner 基，而且已经约化。所有等式在任意系数域中仍成立；若进一步数几何点，才须另看 \(y^3-1\) 是否有重根。

\ProseParagraphBreak
换成字典序 \(y\succ x\)，同一理想的约化基为
\[
\{y-x^2,\ x^3-1\}.
\]
因为 \(y-x^2=-f_1\)、\(x^3-1=xf_1+f_2\)，反向有 \(f_2=x(y-x^2)+(x^3-1)\)；首单项式为 \(y,x^3\)，互素，乘积判据给出 Gröbner 性。两种次序分别先消去 \(x\) 与 \(y\)，商代数虽然相同，适合展示的坐标却不同。

\ProseParagraphBreak
等式\eqref{b3:eq:groebner-cubic-certificates}还有一个实际作用：检查者只需展开有限个多项式等式，再检查一个 S 多项式，就能确认最终答案。只报告“软件返回这组基”，或只验证原多项式对新组的余式为零，都少了一半理想包含的依据。
